\chapter{Introducción}

\section{Motivación y problema profesional}
La digitalización permite consultar y preservar colecciones musicales, pero una resolución
insuficiente puede degradar líneas de pentagrama, cabezas y plicas de nota, alteraciones, texto y
otros elementos pequeños. La superresolución puede mejorar la legibilidad aparente de estas
imágenes, aunque una reconstrucción visualmente convincente no garantiza que conserve la
información musical. El problema profesional de este TFG es, por tanto, determinar cuándo y cómo
puede utilizarse esta técnica sin introducir cambios plausibles pero musicalmente significativos.

Los destinatarios principales son profesionales de digitalización, bibliotecas, archivos,
conservación y patrimonio musical. Como usuarios secundarios se consideran intérpretes,
investigadores y editoriales. Para todos ellos, la utilidad de una técnica dependerá no solo de la
calidad visual, sino también de la fidelidad de la notación, el coste computacional, las licencias
y las condiciones concretas de uso.

La elección del problema combina interés formativo y relevancia aplicada. Permite integrar visión
por computador, aprendizaje profundo, tratamiento de datos y evaluación estadística en un objeto
con estructura semántica exigente. A diferencia de una fotografía natural, una variación mínima
puede cambiar una alteración, un dígito o la continuidad de una línea. Esta propiedad convierte la
partitura en un caso especialmente útil para aprender a distinguir mejora visual, fidelidad y
seguridad de uso. El trabajo cuenta con la orientación conjunta de Elena Vázquez Barrachina y
Jorge Calvo Zaragoza; el diseño, la implementación, la revisión de resultados y las conclusiones
son responsabilidad del autor.

\section{Objetivo general y objetivos específicos}
El objetivo general es evaluar un conjunto pequeño y científicamente contrastado de técnicas de
superresolución sobre degradaciones controladas de partituras digitalizadas, medir su fidelidad,
coste y errores de notación, y convertir la evidencia en un marco de decisión por método,
degradación y uso profesional. Una conclusión válida puede ser que ningún método resulte adecuado
para una condición determinada.

Los objetivos específicos son:
\begin{itemize}
  \item construir un estado del arte trazable basado en fuentes primarias y oficiales;
  \item auditar e inventariar de forma inmutable el conjunto SMB y preservar una muestra
        confirmatoria intacta para la comparación principal;
  \item predeclarar degradaciones, comparadores, unidades independientes, resultados y límites de
        interpretación antes de observar resultados finales;
  \item comparar los métodos sobre entradas emparejadas mediante evidencia de fidelidad, recursos
        y fallos específicos de notación;
  \item evaluar una adaptación acotada con obras separadas de la muestra confirmatoria y sin
        consultar la prueba durante el entrenamiento;
  \item comprobar, sin reajuste posterior, la transferencia del modelo adaptado a una colección
        externa mediante un piloto profesional y un demostrador local reversible; y
  \item producir recomendaciones profesionales reproducibles y limitadas a la evidencia obtenida.
\end{itemize}

Estos objetivos se concretan en seis preguntas de evaluación:
\begin{enumerate}
  \item ¿qué diferencia de PSNR-Y y SSIM-Y presenta cada red respecto de bicúbica dentro de cada
        combinación predeclarada de escala y severidad?;
  \item ¿cambia la relación entre EDSR y SwinIR al aumentar la escala o la severidad?;
  \item ¿qué defectos de pentagramas, símbolos, texto o textura aparecen en los casos cualitativos
        fijados y cuál es su origen observable: entrada degradada, reconstrucción o falta de
        recuperación?;
  \item ¿qué compromiso existe entre fidelidad, tiempo de inferencia, licencias y riesgo de uso?; y
  \item ¿mejora un ajuste fino acotado de EDSR la fidelidad sobre obras SMB no vistas, manteniendo
        separadas las funciones de entrenamiento, validación y prueba?; y
  \item ¿mantiene ese modelo su ventaja al aplicarse, sin reajuste, a obras de otra colección bajo
        las mismas degradaciones sintéticas, y resultan sus salidas utilizables como copias de
        consulta?
\end{enumerate}
No se formuló una hipótesis nula única ni se aplicó un contraste confirmatorio global. Los
intervalos emparejados cuantifican incertidumbre descriptiva sobre comparaciones predeclaradas;
la relevancia profesional requiere además magnitud, evidencia cualitativa y coste.

\section{Impacto esperado y alcance}
La contribución obtenida es una guía de decisión que identifica condiciones de uso, límites y
casos en los que no debe recomendarse la superresolución. Incluye una comparación principal de
modelos preentrenados, un estudio secundario de adaptación dentro de SMB y, tras asegurar ambos, un
piloto sobre doce obras de una colección externa acompañado de un demostrador local de imágenes.
Las conclusiones quedan restringidas a las poblaciones, métodos y degradaciones controladas
realmente evaluadas. El piloto permite estudiar transferencia entre colecciones bajo LR sintética,
pero no demuestra restauración de entradas de baja resolución adquiridas naturalmente. Quedan
fuera la reconstrucción de PDF, el reconocimiento óptico de música y el despliegue público.

\section{Metodología}
El estudio siguió un enfoque de contrato primero. Antes de la evaluación final se fijaron la identidad
del conjunto de datos, las unidades independientes, el manifiesto, los operadores de degradación,
los comparadores, las métricas, el muestreo cualitativo y las reglas de interpretación. La muestra
principal SMB no se empleó para elegir modelos, puntos de control, hiperparámetros ni umbrales. El
estudio secundario congeló antes de entrenar roles por obra sin solapamiento y excluyó las 64 obras
de la comparación principal. El piloto externo fijó después sus doce obras, seis condiciones,
comparadores y casos de revisión antes de generar ninguna salida. Las páginas reciben entradas
idénticas para todos los comparadores y se agrupan por obra o fuente auditada, mientras que las
regiones se tratan como observaciones anidadas.

La evaluación combinó medidas de fidelidad con una revisión sistemática de discontinuidades de
pentagrama, pérdida o deformación de símbolos, uniones indebidas, alteraciones de texto y marcas
alucinadas. También se registró el tiempo de inferencia y el entorno de ejecución. Cada afirmación
material se vinculó a evidencia revisada, con denominador y limitaciones explícitos.

\section{Estructura de la memoria}
Tras esta introducción, el capítulo 2 sintetiza el estado del arte y delimita la brecha; el capítulo
3 analiza alternativas, restricciones y riesgos; y el capítulo 4 documenta los datos y su gobierno.
Los capítulos posteriores describen metodología, resultados, validación profesional, conclusiones
y trabajo futuro. Los anexos conservarán la información de reproducibilidad y, al final, la
reflexión sobre los Objetivos de Desarrollo Sostenible compartida con su entrega independiente.
